% ======================================================================
% Ontology of Continua -- Core 1.3.3
% Cross-K module: crossk_k9_k10.tex
% Cross-level relations between K9 and K10
% ======================================================================

\subsubsection{Межуровневая структура K9–K10}
\label{sec:crossk-k9-k10}

Переход $K_9 \rightarrow K_{10}$ обозначает возникновение
\textbf{метатеоретического континуума} из теоретического континуума. Если
$K_9$ содержит научные теории, парадигмы, модели и логические каркасы, то
$K_{10}$ вводит:
\begin{itemize}
  \item метаоси для сравнения и преобразования целых теоретических систем,
  \item функториальные отображения между категориями моделей,
  \item пороги перекрестной совместимости теорий и метасогласованности,
  \item потоки интерпретации, трансляции и реконструкции,
  \item метациклы, регулирующие эволюцию теорий,
  \item рекурсивные структуры, способные описывать $K_0 \dots K_{10}$,
        включая самих себя.
\end{itemize}

$K_{10}$ — наиболее высокая репрезентативная ступень в модели Core, которая
ещё находится в пределах эпистемического контура человека.

% ----------------------------------------------------------------------
\subsubsection{1. Задействованные уровни и их роли}

\paragraph{K9 (теоретический континуум).}
$K_9$ предоставляет:
\begin{itemize}
  \item логические, онтологические и методологические оси $A^9$,
  \item теоретические потенциалы $P^9$ (объяснительная сила, когерентность,
        простота),
  \item пороги $\Theta^9$ (согласованность, эмпирическая адекватность),
  \item потоки $J^9$ (аргументы, доказательства, критика),
  \item парадигменные циклы $C^9$ (нормальная наука, аномалия, революция).
\end{itemize}

Однако $K_9$ не включает:
\begin{itemize}
  \item структурированное пространство сравнения теорий,
  \item функториальные отображения между модельными структурами,
  \item мета-логику и мета-онтологию,
  \item рекурсивное представление, включающее теорию теорий.
\end{itemize}

\paragraph{K10 (метатеоретический континуум).}
$K_{10}$ вводит:
\begin{itemize}
  \item оси $A^{10}$ (типы функторов, мета-логика, онтология онтологий),
  \item потенциалы $P^{10}$ (метасогласованность, экспрессивность, универсальность),
  \item пороги $\Theta^{10}$ (категориальная согласованность, функториальная
        совместимость, мета-логические пределы),
  \item потоки $J^{10}$ (интерпретация, трансляция, реконструкция,
        функториальные подъёмы),
  \item циклы $C^{10}$ (мета-стабильность, само-согласованность,
        перекрестно-парадигмальное согласование),
  \item меру $k(K_{10})$ глобальной мета-согласованности.
\end{itemize}

Определяющая характеристика $K_{10}$ — его рекурсивная способность:
\[
K_{10} \text{ описывает все уровни } K_0\dots K_{10}.
\]

% ----------------------------------------------------------------------
\subsubsection{2. Совместные и наследуемые оси и пороги}

\paragraph{Наследованная структура.}
Теории в $K_9$ дают объекты; $K_{10}$ даёт морфизмы:
\[
\text{Models}(K_9) \rightarrow \text{Functors}(K_{10}).
\]

Символические структуры, логика и методологии из $K_9$ становятся элементами в
категориальной архитектуре $K_{10}$.

\paragraph{Новые оси в \texorpdfstring{$K_{10}$}{K_{10}}.}
К ним относятся:
\begin{itemize}
  \item $A_{\mathrm{functor}}$ — типы преобразований моделей,
  \item $A_{\mathrm{meta}}$ — мета-логические и мета-онтологические различия,
  \item $A_{\mathrm{compat}}$ — степени перекрестной совместимости теорий,
  \item $A_{\mathrm{reflect}}$ — оси для самоописания.
\end{itemize}

\paragraph{Мета-пороги.}
$\Theta^{10}$ содержит:
\begin{itemize}
  \item \textbf{порог согласованности} — диаграммы категорий должны
        коммутировать,
  \item \textbf{порог совместимости} — функторы должны сохранять структуру между
        теориями,
  \item \textbf{мета-логический порог} — исключение парадокса или
        неконсистентности,
  \item \textbf{порог экспрессивности} — минимальная способность представлять
        различия на всех нижележащих уровнях.
\end{itemize}

Нарушение $\Theta^9$ препятствует формированию $K_{10}$.

% ----------------------------------------------------------------------
\subsubsection{3. Межуровневые потоки, циклы и напряжения}

\paragraph{Потоки.}
Потоки в $K_{10}$ расширяют теоретические потоки:
\[
J^{10} = (J_{\mathrm{interpret}},\ J_{\mathrm{translate}},\
         J_{\mathrm{reconstruct}},\ J_{\mathrm{lift}}),
\]
где:
\begin{itemize}
  \item $J_{\mathrm{interpret}}$ — переинтерпретация теории в новой рамке,
  \item $J_{\mathrm{translate}}$ — отображения между формализмами,
  \item $J_{\mathrm{reconstruct}}$ — восстановление теоретической структуры,
  \item $J_{\mathrm{lift}}$ — представление теории как элемента более высокой
        категории или метамодели.
\end{itemize}

\paragraph{Циклы.}
Метациклы $C^{10}$ включают:
\[
C_{\mathrm{meta}}:\
\text{model} \rightarrow \text{interpretation} \rightarrow \text{translation}
\rightarrow \text{refinement} \rightarrow \text{model}.
\]

Эти циклы обеспечивают стабильность мета-континуума.

\paragraph{Напряжение.}
Межуровневое напряжение:
\[
T_{9,10} =
f(\Theta^{10} - P^{10},\ A^{10} - A^9,\
  J^{10} - J^9,\ \Theta^9 - \Theta^{10}).
\]
Если $T_{9,10}$ превышает размерностный порог, $K_{10}$ коллапсирует в
некатегориальный теоретический режим.

% ----------------------------------------------------------------------
\subsubsection{4. Условия рождения и исчезновения между уровнями}

\paragraph{Рождение $K_{10}$.}
Метатеоретический континуум возникает, когда:
\[
T_9 > \Theta_{9,\mathrm{dim}}
\quad\text{и}\quad
A^{10} \in M_{10} \setminus A(K_9).
\]

Расширение:
\[
\Omega(K_{10}) = \Omega(K_9) \cup \Delta\Omega_{\mathrm{meta}}.
\]

\paragraph{Распространение гибели.}
Если теоретическая устойчивость нарушена ($\Theta^9$), $K_{10}$ не может сформироваться.

Если нарушены мета-пороги $\Theta^{10}$ (несовместимость функторов,
мета-логический парадокс), $K_{10}$ коллапсирует, но $K_9$ остаётся.

% ----------------------------------------------------------------------
\subsubsection{5. Концептуальные примеры}

\paragraph{Сравнение теорий.}
$K_{10}$ позволяет структурированно сравнивать теории через функторы:
\[
\mathcal{F} : \mathsf{Mod}(T_1) \rightarrow \mathsf{Mod}(T_2).
\]

\paragraph{Интерпретация парадигм.}
Несовместимые теории в $K_9$ могут быть согласованы путём подъёма их в
мета-рамку $K_{10}$.

\paragraph{Избежание парадокса.}
Если мета-рамка нарушает мета-логический порог
$\Theta^{10}_{\mathrm{coh}}$, разрушается весь мета-континуум.

Следовательно, переход K9→K10 в сжатой форме — это
появление мета-слоя, способного сравнивать, преобразовывать и реконструировать
целые теоретические системы, включая себя самого.
